\chapter{Trésorerie et financement}\label{ch:fm:treasury-and-funding}

En mars 2020, les exigences de marge initiale auprès des chambres de compensation mondiales ont augmenté d'environ 300~milliards de dollars en un mois, et la marge de variation appelée
dans les marchés compensés comme dans les marchés bilatéraux a été bien plus forte qu'en février. Toute société ayant des positions compensées ou soumises à marge devait trouver les liquidités le
matin où chaque appel tombait, là où sa trésorerie l'avait laissée~: chez un courtier qui acceptait de la libérer, dans un fonds monétaire pouvant
être vendu le jour même, ou sur une ligne de crédit que la banque honorerait. Une société qui avait gardé sa trésorerie là où elle ne pouvait pas l'atteindre
en a appris à ses dépens la portée réelle de son \omterm{def:fm:treasury-and-funding:horizon}{horizon de survie}.

\section{Ce que fait la trésorerie dans une société de trading}

La trésorerie d'une société de trading détient et déplace ses liquidités et son collatéral. Elle garde assez de fonds chez chaque prime broker, courtier compensateur et chambre de compensation pour
satisfaire leurs exigences~; elle finance les positions longues et emprunte les titres des positions courtes au moindre coût~; elle prévoit ce que les prochains jours
exigeront~; elle garde une réserve pour les jours où la prévision se trompe. C'est la fonction dont la défaillance est la plus rapide~: une société aux bonnes positions et
sans liquidités pour un appel de marge est en défaut dès le soir.

\begin{definition}[Trésorerie non grevée, coussin de liquidité]\label{def:fm:treasury-and-funding:cash}
La \emph{trésorerie non grevée}\index{trésorerie non grevée} (unencumbered cash) est la trésorerie, ainsi que les actifs convertibles en liquidités le jour même, que la société détient libres de tout nantissement,
de toute exigence de marge ou restriction et peut les utiliser sur-le-champ. Un \emph{coussin de liquidité}\index{coussin de liquidité} (liquidity buffer) est la part qui en est réservée contre les sorties en situation de
tension (appels de marge, restitutions de collatéral, rachats, perte de financement) sur un horizon défini, et non utilisée pour financer des positions.
\end{definition}

\begin{omfigure}
\begin{tikzpicture}[font=\footnotesize,>=stealth,box/.style={draw,rounded corners,align=center,minimum height=0.75cm,text width=2.4cm}]
\node[box,fill=omDef!15,text width=3cm] (t) at (0,0) {trésorerie\\{\scriptsize trésorerie non grevée}};
\node[box,fill=omProp!15] (a) at (-4.2,1.9) {prime broker A};
\node[box,fill=omProp!15] (b) at (0,2.3) {prime broker B};
\node[box,fill=omProp!15] (c) at (4.2,1.9) {prime broker C};
\node[box,fill=omThm!15] (f) at (4.6,-2.3) {courtier compensateur\\{\scriptsize future}};
\node[box,fill=gray!15] (ccp) at (4.6,-3.8) {chambre de compensation};
\node[box,fill=gray!15] (mmf) at (-4.6,-2.3) {fonds monétaire, bons du Trésor};
\node[box,fill=gray!15] (rcf) at (0,-2.8) {ligne de crédit confirmée};
\draw[<->] (t) -- (a); \node[font=\scriptsize] at (-2.9,0.75) {appels, restitutions}; \draw[<->] (t) -- (b); \draw[<->] (t) -- (c);
\draw[<->] (t) -- node[pos=0.5,above,sloped,font=\scriptsize] {marge de variation} (f); \draw[<->] (f) -- (ccp);
\draw[<->] (t) -- node[pos=0.5,above,sloped,font=\scriptsize] {vente le jour même} (mmf); \draw[->] (rcf) -- node[right,font=\scriptsize] {tirage} (t);
\end{tikzpicture}

\omcaption{La trésorerie d'une société de trading~: les comptes qui détiennent son collatéral, appellent ou restituent la marge chaque jour, et les sources qu'elle peut convertir en liquidités
le jour même. Chaque flèche a ses propres conditions, heures limites et délais de préavis. Schéma de principe.}\label{fig:fm:treasury-and-funding:map}
\end{omfigure}

La version bancaire de l'idée est le ratio de liquidité à court terme (LCR) du livre~6, et le prix interne de l'usage du bilan est son taux de cession interne (funds-transfer pricing, FTP).
Une société pour compte propre ou un fonds n'a, dans la plupart des juridictions, aucun ratio réglementaire à respecter~; elle dispose des conditions de ses courtiers, de ses propres positions et des mouvements du marché.

\section{Relations avec les prime brokers et conditions de marge}

Un prime broker (livre~1) finance les positions longues de la société, lui prête les titres qu'elle vend à découvert, détient ses positions et son collatéral, et fixe la marge qu'il exige
sur elles. Le minimum réglementaire des comptes clients américains n'est que le plancher (\cref{dat:fm:treasury-and-funding:rules})~; ce qui contraint un
client professionnel est la règle propre du courtier.

\begin{definition}[Marge maison]\label{def:fm:treasury-and-funding:house}
La \emph{marge maison}\index{marge maison} (house margin) est la marge qu'un courtier exige d'un client selon ses propres règles, au-dessus de tout minimum réglementaire ou imposé par la
chambre de compensation~: typiquement un taux sur l'exposition brute, un taux sur l'exposition nette, des majorations pour concentration, pour liquidité et pour corrélation, et une méthodologie que le courtier
peut modifier à sa discrétion, sauf si le contrat en dispose autrement.
\end{definition}

\begin{definition}[Immobilisation de marge]\label{def:fm:treasury-and-funding:lockup}
Une \emph{immobilisation de marge}\index{immobilisation de marge} (margin lock-up) est l'engagement contractuel d'un courtier à ne pas modifier ses conditions de marge (taux, méthodologie, collatéral
éligible) ni retirer son financement à un client sans délai de préavis, en général de trente à quatre-vingt-dix jours, sauf lors d'événements définis, comme une
forte baisse de la valeur liquidative du client.
\end{definition}

\begin{dated}{2026-09}{Planchers réglementaires de marge pour les comptes clients américains}\label{dat:fm:treasury-and-funding:rules}
La Regulation T (12 CFR 220.12) fixe la marge initiale sur un titre de capital acheté sur marge à 50~\% de sa valeur de marché courante, et sur une vente à découvert d'un
titre non exempté à 150~\% de sa valeur. La FINRA Rule 4210 fixe une marge de maintenance de 25~\% de la valeur de marché courante des titres longs sur marge,
et permet aux sociétés membres d'appliquer une marge de portefeuille fondée sur le risque au lieu des règles par stratégie. Les \omterm{def:fm:treasury-and-funding:house}{marges maison} des courtiers se situent au-dessus de ces planchers~; les
accords offshore et institutionnels sont fixés par contrat.
\end{dated}

Trois courtiers aux règles différentes valorisent très différemment le même portefeuille. Le portefeuille du chapitre est long de 390~millions de dollars et court de 330~millions de dollars sur dix
positions en actions, couvert par un future d'indice vendu de 60~millions de dollars chez un courtier compensateur (illustratif). Le courtier A facture 4~\% du brut plus 25~\% de
la valeur absolue de l'exposition nette, avec une majoration de 15~\% sur toute position au-delà de 60~millions de dollars~; le courtier B facture un taux forfaitaire de 10~\% du brut~; le courtier C facture 5~\% du brut plus
15~\% du net, avec une majoration de 30~\% au-delà de 40~millions de dollars. Une seule position longue de 120~millions de dollars coûte 43.8~millions de dollars de marge chez A, 12.0~millions de dollars chez B et 48.0
millions de dollars chez C~: A et C récompensent un portefeuille équilibré et pénalisent la concentration, B fait payer la seule taille.

\omcode{code/firm/treasury/firm_treasury.py}{35}{55}{Règles de marge maison sous forme de données~: termes sur le brut, sur le net et pour la concentration, et spreads de financement sur les
positions longues et courtes.}\label{lst:fm:treasury-and-funding:rules}

\section{Optimisation des marges}

\begin{definition}[Optimisation des marges]\label{def:fm:treasury-and-funding:optim}
L'\emph{optimisation des marges}\index{optimisation des marges} (margin optimisation) est le choix du lieu où détenir chaque position et de la manière d'apparier les positions dans chaque compte,
de sorte que le total de la marge et le coût de financement sur l'ensemble des courtiers et chambres de compensation d'une société soit le plus faible pour le même portefeuille.
\end{definition}

\begin{proposition}[Optimiser ne coûte jamais plus cher que le recours à un seul courtier]\label{prop:fm:treasury-and-funding:lp}
Si la marge de chaque courtier est une somme de taux multipliés par l'exposition brute, l'exposition nette en valeur absolue et les excédents des positions par rapport aux caps, et si le financement est linéaire
dans les positions, le coût annuel est une fonction convexe, linéaire par morceaux, des fractions $x_{ib}\ge0$ de chaque position $i$ détenue chez le courtier $b$,
$\sum_bx_{ib}=1$. Son minimum est la solution d'un programme linéaire, et il n'est pas supérieur au coût de détenir l'ensemble du portefeuille chez un seul
courtier.
\end{proposition}

\begin{proof}
L'exposition brute est linéaire en $x$~; la valeur absolue d'une fonction linéaire et la partie positive d'une fonction affine sont convexes~; les sommes de fonctions convexes
à coefficients positifs ou nuls sont convexes. Remplacer chaque valeur absolue par une variable $t_b\ge\pm\sum_ip_ix_{ib}$ et chaque partie positive par
$u_{ib}\ge|p_i|x_{ib}-\text{cap}_b$, $u_{ib}\ge0$, donne un programme linéaire de même minimum. Tout détenir chez un seul courtier est un
point admissible, donc le minimum n'est pas supérieur à son coût.
\end{proof}

\begin{table}[ht]
\centering\small
\begin{tabular}{lrrrrr}
\toprule
allocation & marge A & marge B & marge C & marge totale & coût annuel \\
\midrule
tout chez A & 71.55 & -- & -- & 71.55 & 6.88 \\
tout chez B & -- & 72.00 & -- & 72.00 & 5.99 \\
tout chez C & -- & -- & 144.00 & 144.00 & 10.92 \\
au prorata & 14.60 & 24.00 & 15.00 & 53.60 & 5.82 \\
optimisée & 20.40 & 6.00 & 7.50 & 33.90 & 4.98 \\
\bottomrule
\end{tabular}
\caption{\omterm{def:fm:treasury-and-funding:house}{Marge maison} et coût annuel (marge, à un coût de trésorerie de 5~\%, plus le financement) du portefeuille du chapitre selon cinq allocations réparties entre trois
prime brokers, en millions de dollars. Données~: \texttt{fm\_treasury.summary}.}\label{tab:fm:treasury-and-funding:alloc}
\end{table}

L'optimiseur équilibre la tranche de chaque courtier, de sorte que la charge nette du courtier A et les majorations pour concentration de C disparaissent presque, et n'envoie à B que
ce que son taux brut forfaitaire tarifie à bon marché (le \cref{tab:fm:treasury-and-funding:alloc}, la \cref{fig:fm:treasury-and-funding:margins}). La marge passe de 72.0~millions de dollars
chez le moins cher des courtiers pris isolément à 33.9~millions de dollars, soit 53~\% de moins, et le coût annuel passe de 5.99~millions de dollars à 4.98~millions de dollars. Répartir le livre
à parts égales ne fait qu'environ la moitié du chemin, jusqu'à 53.6~millions de dollars.

\begin{omfigure}
\begin{tikzpicture}
\begin{axis}[width=10.5cm,height=5cm,ybar stacked,bar width=18pt,ylabel={\small marge maison (millions de dollars)},ymin=0,ymax=150,xtick={0,1,2,3,4},
  xticklabels={tout chez A,tout chez B,tout chez C,au prorata,optimisée},xticklabel style={text width=1.6cm,align=center},xmin=-0.6,xmax=4.6,
  tick label style={font=\footnotesize},grid=major,grid style={gray!20},table/col sep=comma,
  legend cell align=left,legend style={font=\scriptsize,at={(0.5,-0.25)},anchor=north,/tikz/every even column/.append style={column sep=8pt}},legend columns=3]
\addplot[fill=omDef!55,draw=omDef] table[x=pos,y=A]{figdata/desk/14-treasury-and-funding/margins.csv};
\addplot[fill=omProp!55,draw=omProp] table[x=pos,y=B]{figdata/desk/14-treasury-and-funding/margins.csv};
\addplot[fill=omThm!50,draw=omThm] table[x=pos,y=C]{figdata/desk/14-treasury-and-funding/margins.csv};
\legend{courtier A,courtier B,courtier C}
\end{axis}
\end{tikzpicture}

\omcaption{\omterm{def:fm:treasury-and-funding:house}{Marge maison} du même livre selon cinq allocations réparties entre trois prime brokers. Le programme linéaire de
la \cref{prop:fm:treasury-and-funding:lp} équilibre la tranche de chaque courtier et réduit la marge de 53~\% par rapport au courtier le moins cher pris isolément. Données~:
\texttt{fm\_treasury.summary}.}\label{fig:fm:treasury-and-funding:margins}
\end{omfigure}

\begin{remark}[Ce que l'optimiseur ne voit pas]\label{rem:fm:treasury-and-funding:limits}
Les courtiers tarifient des relations, non des tranches~: un courtier qui ne voit que le reliquat équilibré d'un livre peut le retarifer, et celui qui n'en détient que peu
a peu de raisons de soutenir la firme en stress. Déplacer des positions entre courtiers coûte de l'argent et du temps, et en stress cela peut être impossible.
L'optimiseur donne l'allocation la moins chère pour les conditions d'aujourd'hui~; le trésorier en retient une qui reste acceptable quand les conditions changent.
\end{remark}

\section{Prévoir la trésorerie et le collatéral}

Un prévisionnel de trésorerie suit chaque compte sur les jours à venir. Chaque compte dispose de collatéral~; chaque jour il gagne ou perd le P\&L des positions
qu'il détient et fait face à une exigence~; un déficit est un appel réglé sur la \omterm{def:fm:treasury-and-funding:cash}{trésorerie non grevée} le jour même, et un excédent retourne à la trésorerie si le courtier
la libère (\cref{lst:fm:treasury-and-funding:forecast}). L'exigence du courtier compensateur sur le contrat à terme provient de la simulation historique filtrée de firm.initmargin
(livre~6), et celle des prime brokers découle de leurs règles maison multipliées par un coefficient qu'ils peuvent relever en période de stress.

\omcode{code/firm/treasury/firm_treasury.py}{111}{134}{Le prévisionnel de trésorerie~: le P\&L de chaque compte, l'exigence, l'appel ou la restitution quotidiens, et
l'horizon de survie.}\label{lst:fm:treasury-and-funding:forecast}

\begin{definition}[Horizon de survie]\label{def:fm:treasury-and-funding:horizon}
Le \emph{horizon de survie}\index{horizon de survie} d'une firme dans un scénario est le nombre de jours au bout duquel sa \omterm{def:fm:treasury-and-funding:cash}{trésorerie non grevée}, après avoir honoré
chaque appel de marge et chaque décaissement contractuel qu'engendre le scénario, passe pour la première fois sous zéro, sans nouveau financement.
\end{definition}

\section{Tutoriel~: trois courtiers et une mauvaise semaine}\label{tut:fm:treasury-and-funding:tut}

\begin{tutorial}
\textbf{Objectif.} Répartir le livre entre trois courtiers, puis prévoir trente jours de trésorerie dans un scénario de référence et dans une semaine de stress, avec et sans
\omterm{def:fm:treasury-and-funding:lockup}{immobilisation de marge}. \textbf{État final~:} le \cref{tab:fm:treasury-and-funding:alloc} et le graphique de la piste de trésorerie (la \cref{fig:fm:treasury-and-funding:runway}).
\begin{tutsteps}
\item \textbf{L'allocation.} \texttt{fm\_treasury.allocations()} construit les cinq allocations~; \texttt{firm.treasury.optimise} résout le
programme linéaire avec le solveur HiGHS de scipy.
\item \textbf{Les comptes.} Chaque compte chez un prime broker démarre avec sa marge plus un coussin de 10~\%~; la firme dispose de 40~millions de dollars de \omterm{def:fm:treasury-and-funding:cash}{trésorerie non grevée}.
\item \textbf{Le stress.} Pendant cinq jours, les positions longues baissent de 2~\% par jour, les shorts de 1~\% et l'indice de 2~\%~; suivent cinq jours plats puis une reprise lente. Les
courtiers relèvent leur \omterm{def:fm:treasury-and-funding:house}{marge maison} d'un cinquième par jour, jusqu'à la doubler au jour 5, la maintiennent jusqu'au jour 15, puis la ramènent à une fois et demie.
\item \textbf{La piste.} \texttt{fm\_treasury.run(stress, notice)} prévoit la trésorerie avec une immobilisation de \texttt{notice} jours (0 pour aucune).
\end{tutsteps}
\end{tutorial}

\begin{omfigure}
\begin{tikzpicture}
\begin{axis}[width=10.5cm,height=5.4cm,xlabel={\small jour},ylabel={\small trésorerie non grevée (millions de dollars)},xmin=0,xmax=30,ymin=-35,ymax=50,
  tick label style={font=\footnotesize},grid=major,grid style={gray!20},table/col sep=comma,
  legend cell align=left,legend style={font=\scriptsize,at={(0.5,-0.3)},anchor=north,/tikz/every even column/.append style={column sep=8pt}},legend columns=2]
\addplot[gray,thick] table[x=day,y=base]{figdata/desk/14-treasury-and-funding/runway.csv};
\addplot[omThm,thick,mark=*,mark size=1pt] table[x=day,y=lockup]{figdata/desk/14-treasury-and-funding/runway.csv};
\addplot[omDef,very thick] table[x=day,y=stress]{figdata/desk/14-treasury-and-funding/runway.csv};
\addplot[omProp,thick,dashed] table[x=day,y=one]{figdata/desk/14-treasury-and-funding/runway.csv};
\draw[black,thin] (axis cs:0,0) -- (axis cs:30,0);
\legend{cas de référence,{stress, immobilisation de 30 jours},{stress, sans immobilisation},{stress, tout chez le courtier B}}
\end{axis}
\end{tikzpicture}

\omcaption{Trente jours de \omterm{def:fm:treasury-and-funding:cash}{trésorerie non grevée}. Sans immobilisation, le livre optimisé est à court de trésorerie dès le jour 5, et le même livre chez un seul courtier dès le
jour 3~; avec une immobilisation de trente jours, il conserve 26.7~millions de dollars au pire moment. Données~: \texttt{fm\_treasury.run}.}\label{fig:fm:treasury-and-funding:runway}
\end{omfigure}

La semaine de stress coûte au livre 15.5~millions de dollars de P\&L~; les hausses de marge des courtiers coûtent davantage. Sans immobilisation, l'exigence totale passe de
43.0~millions de dollars au jour 1 à 71.8~millions de dollars au jour 5 (\cref{fig:fm:treasury-and-funding:req}), la firme verse 57.9~millions de dollars d'appels de marge sur le mois,
et sa trésorerie de 40~millions de dollars est épuisée dès le jour 5~: il lui aurait fallu un \omterm{def:fm:treasury-and-funding:cash}{coussin de liquidité} supérieur de 8.1~millions de dollars. Le même livre chez le seul courtier B est épuisé dès le jour 3 et aurait nécessité
29.8~millions de dollars de plus. Avec une immobilisation de trente jours, l'exigence reste proche de 37~millions de dollars, les appels totalisent 23.0~millions de dollars, et la firme termine le pire jour avec
26.7~millions de dollars. Une immobilisation plus courte que le stress ne fait que reporter la rupture~: cinq jours de préavis la reportent au jour 6, dix jours au jour 11. La
marge du courtier sur le contrat à terme bouge à peine, de 2.02~millions de dollars à 2.13~millions de dollars, car son modèle réagit à la volatilité avec un retard~; la part
importante tient à la discrétion des courtiers.

\begin{omfigure}
\begin{tikzpicture}
\begin{axis}[width=10.5cm,height=4.8cm,xlabel={\small jour},ylabel={\small exigence totale (millions de dollars)},xmin=1,xmax=30,ymin=0,ymax=80,
  tick label style={font=\footnotesize},grid=major,grid style={gray!20},table/col sep=comma,
  legend cell align=left,legend style={font=\scriptsize,at={(0.5,-0.32)},anchor=north,/tikz/every even column/.append style={column sep=8pt}},legend columns=2]
\addplot[omDef,very thick,const plot] table[x=day,y=nolock]{figdata/desk/14-treasury-and-funding/req.csv};
\addplot[omThm,thick,dashed,const plot] table[x=day,y=lock]{figdata/desk/14-treasury-and-funding/req.csv};
\legend{sans immobilisation,immobilisation de 30 jours}
\end{axis}
\end{tikzpicture}

\omcaption{L'exigence de marge totale sur les quatre comptes pendant le mois de stress. L'écart entre les courbes mesure la hausse de la \omterm{def:fm:treasury-and-funding:house}{marge maison} décidée par les courtiers, que
une immobilisation reporte. Données~: \texttt{fm\_treasury.paths}.}\label{fig:fm:treasury-and-funding:req}
\end{omfigure}

\section{Financement en période de stress}

Le stress actionne simultanément tous les leviers du prévisionnel~: les positions perdent de l'argent, les exigences montent, les courtiers cessent de restituer l'excédent, et les sources de financement
qui étaient bon marché deviennent indisponibles. La spirale des marges du livre~3 en est la version de marché~: les ventes forcées pour couvrir les appels font bouger les prix, ce qui alourdit
les appels ailleurs. Les défenses de la firme sont le \omterm{def:fm:treasury-and-funding:cash}{coussin de liquidité}, les conditions et la diversité de ses financements.

La diversité a trois volets. Plusieurs courtiers, de sorte qu'une décision d'un seul courtier ne déplace qu'une partie de l'exigence~; plusieurs sources de trésorerie, afin
qu'un fonds monétaire qui gèle les rachats ou une banque qui refuse de prêter ne soit pas la seule voie~; et des actifs dont la valeur ne baisse pas avec celle du livre, puisque le
collatéral que la firme apporte en stress vaut ce qu'il trouvera preneur ce jour-là. Une ligne de crédit confirmée est la plus coûteuse des trois en période calme,
car la firme paie une commission sur les fonds qu'elle n'utilise pas, et c'est celle dont la valeur est la plus haute le jour où arrive un appel auquel
rien d'autre ne peut faire face (\cref{fig:fm:treasury-and-funding:map}).

\begin{dated}{2026-09}{Recommandations internationales sur la préparation aux marges}\label{dat:fm:treasury-and-funding:fsb}
Après l'épisode de mars 2020, le BCBS, le CPMI et l'OICV ont réexaminé les pratiques de marge (rapport final, septembre 2022), et le Conseil de stabilité financière
a publié huit recommandations à l'intention des intervenants de marché non bancaires sur la préparation de la liquidité face aux appels de marge et de collatéral (rapport final, 10
décembre 2024)~: intégrer le risque de liquidité des appels de marge et de collatéral à leur gestion des risques et à leur gouvernance, avec des plans de financement d'urgence~;
soumettre leur liquidité à des tests de résistance pour de tels appels, dans des scénarios extrêmes mais plausibles~; et détenir des liquidités suffisantes et des actifs liquides
diversifiés et facilement disponibles, avec des dispositifs de collatéral appropriés et des contacts réguliers avec leurs contreparties.
\end{dated}

\begin{method}[Fixer le coussin de liquidité]\label{met:fm:treasury-and-funding:buffer}
\begin{enumerate}
\item Recenser chaque compte~: ce qu'il détient, sa règle d'exigence, si l'excédent est restitué, et le préavis avant que ses conditions ne puissent changer.
\item Construire des scénarios de stress qui font bouger ensemble le résultat et les exigences, y compris les hausses discrétionnaires des courtiers et une contrepartie qui cesse
de restituer l'excédent.
\item Prévoir la trésorerie chaque jour sur l'horizon du préavis le plus long~; le coussin est le déficit cumulé le plus élevé.
\item Négocier des immobilisations pour les conditions qui bougent le plus dans les scénarios, et répartir le financement de sorte qu'aucune décision d'un seul courtier n'épuise le coussin.
\item Communiquer chaque mois au \omterm{def:fm:the-risk-management-function:cro}{comité des risques} l'\omterm{def:fm:treasury-and-funding:horizon}{horizon de survie} sous chaque scénario, ainsi qu'après chaque changement substantiel de conditions ou de positions.
\end{enumerate}
\end{method}

Le livre~1 raconte comment les conditions de marge d'un family office auprès de plusieurs banques, et les réactions différentes de ces banques lorsque ses positions ont baissé, ont conduit à de
lourdes pertes pour certaines d'entre elles. Pour l'emprunteur, la leçon est celle de ce chapitre~: les conditions qui comptent sont celles qui peuvent changer
en une semaine.

\section{Construction~: le moteur de trésorerie}\label{bld:fm:treasury-and-funding:treasury}

\begin{build}
\bfield{Rôle} \omterm{def:fm:treasury-and-funding:house}{Marge maison} traitée comme des données, répartition d'un livre entre courtiers, et prévision quotidienne de la trésorerie et du collatéral avec son
\omterm{def:fm:treasury-and-funding:horizon}{horizon de survie}.
\bfield{Interface} \texttt{firm.treasury}~: \texttt{Broker}, \texttt{margin}, \texttt{financing}, \texttt{annual\_cost}, \texttt{optimise}~;
\texttt{fcm\_margin} (sur \texttt{firm.initmargin}), \texttt{house\_path}, \texttt{forecast}, \texttt{survival\_horizon}, \texttt{buffer\_needed}.
\bfield{Règles} Les allocations sont des fractions dont la somme vaut un par position~; le coût de l'optimiseur ne dépasse jamais celui d'un seul courtier~; un appel est payé le
jour même où il survient~; une immobilisation maintient le multiplicateur maison à un pendant son préavis.
\bfield{Tests d'acceptation} \texttt{code/firm/treasury/tests/}~: marge et financement sur des chiffres connus~; l'optimiseur comparé à chaque courtier pris isolément~; une
prévision sur trois jours avec un excédent conservé~; la mise à l'échelle de la marge de compensation.
\bfield{Pour aller plus loin} Éligibilité et décotes du collatéral~; une ligne de crédit renouvelable avec une commission d'engagement~; le stress produit par le moteur de scénarios de firm.margin
pour le courtier compensateur.
\end{build}

\begin{omsources}
\item BCBS, CPMI et IOSCO, \emph{Review of margining practices}, septembre 2022, et communiqué de presse du 29 septembre 2022.
\item Financial Stability Board, \emph{Liquidity Preparedness for Margin and Collateral Calls: Final report}, 10 décembre 2024.
\item 12 CFR 220.12~; FINRA Rule 4210.
\end{omsources}

\section{Exercices}

\begin{exercise}[$\star$]\label{exo:fm:treasury-and-funding:1}
Calculer la \omterm{def:fm:treasury-and-funding:house}{marge maison} d'une seule position longue de 120~millions de dollars chez chacun des trois courtiers du chapitre.
\end{exercise}

\begin{exercise}[$\star$]\label{exo:fm:treasury-and-funding:2}
Quelle est la marge initiale de la Regulation T sur une position longue de 10~millions de dollars dans un titre en marge, et sur une vente à découvert de 10~millions de dollars~?
\end{exercise}

\begin{exercise}[$\star$]\label{exo:fm:treasury-and-funding:3}
La marge du livre passe de 72.0~millions à 33.9~millions de dollars. Au coût de la trésorerie de 5~\%, que vaut cette baisse un an avant financement~?
\end{exercise}

\begin{exercise}[$\star\star$]\label{exo:fm:treasury-and-funding:4}
Montrer que détenir l'ensemble du livre chez un seul courtier est un point réalisable du programme linéaire, et expliquer pourquoi le coût optimisé ne peut pas le dépasser.
\end{exercise}

\begin{exercise}[$\star\star$]\label{exo:fm:treasury-and-funding:5}
Pourquoi la règle du courtier B en fait-elle le courtier unique le moins cher pour ce livre, alors que sa marge n'est pas la plus basse~?
\end{exercise}

\begin{exercise}[$\star\star$]\label{exo:fm:treasury-and-funding:6}
Une immobilisation de dix jours est proposée contre une de trente. Sous le stress du chapitre, qu'achète chacune~?
\end{exercise}

\begin{exercise}[$\star\star\star$]\label{exo:fm:treasury-and-funding:7}
\emph{Programmation.} Relancer le stress avec une trésorerie initiale de 50~et de 60~millions de dollars, pour le livre optimisé et pour le livre chez le courtier B. Quand chacun
est-il à court~?
\end{exercise}

\begin{exercise}[$\star\star\star$]\label{exo:fm:treasury-and-funding:8}
\emph{Trouver la faille.} «~Nous avons tout transféré chez le courtier à la marge la plus basse~; cela libère de la trésorerie, donc notre liquidité est meilleure.~»
\end{exercise}

\section{Problème~: Trois courtiers et une mauvaise semaine}

\begin{problem}[{Problème du week-end --- trois courtiers et une mauvaise semaine}]\label{pb:fm:treasury-and-funding:1}
Un nouveau trésorier doit décider où la firme détient son livre et combien de trésorerie il garde disponible.

\medskip\noindent\textbf{Partie I --- La fonction.}
\begin{enumerate}
\item Définir la \omterm{def:fm:treasury-and-funding:cash}{trésorerie non grevée} et le \omterm{def:fm:treasury-and-funding:cash}{coussin de liquidité}.
\item Que s'est-il passé pour la marge en mars 2020~?
\item Définir la \omterm{def:fm:treasury-and-funding:house}{marge maison} et l'\omterm{def:fm:treasury-and-funding:lockup}{immobilisation de marge}.
\item Énoncer les planchers réglementaires américains pour la marge des clients.
\end{enumerate}

\medskip\noindent\textbf{Partie II --- L'allocation.}
\begin{enumerate}[resume]
\item Définir l'\omterm{def:fm:treasury-and-funding:optim}{optimisation des marges}.
\item Énoncer et démontrer la \cref{prop:fm:treasury-and-funding:lp}.
\item Donner la marge et le coût annuel des cinq allocations.
\item Comment l'optimiseur utilise-t-il la règle de chaque courtier~?
\item Que ne voit pas l'optimiseur~?
\end{enumerate}

\medskip\noindent\textbf{Partie III --- La prévision.}
\begin{enumerate}[resume]
\item Définir l'\omterm{def:fm:treasury-and-funding:horizon}{horizon de survie}.
\item Décrire le scénario de stress et le résultat qu'il provoque.
\item Donner l'exigence totale aux jours 1 et 5, avec et sans immobilisation, ainsi que la marge du courtier de compensation.
\item Donner les \omterm{def:fm:treasury-and-funding:horizon}{horizons de survie} et les coussins nécessaires, pour le livre optimisé et pour le livre chez un seul courtier.
\item Qu'achète une immobilisation plus courte que le stress~?
\end{enumerate}

\medskip\noindent\textbf{Partie IV --- La décision.}
\begin{enumerate}[resume]
\item Que demandent les recommandations internationales de 2024 aux sociétés non bancaires~?
\item Comment fixer le coussin~?
\item Lequel vaut le plus dans ce stress~: l'optimisation ou l'immobilisation~?
\item Que modifier dans les conventions passées avec les courtiers~?
\item Énoncer le \emph{résultat nommé}~: la marge économisée en optimisant sur trois courtiers plutôt qu'en en utilisant un seul, et l'\omterm{def:fm:treasury-and-funding:horizon}{horizon de survie} dans la semaine de stress
avec et sans \omterm{def:fm:treasury-and-funding:lockup}{immobilisation de marge}.
\item En deux phrases, rédiger la politique de trésorerie.
\end{enumerate}
\end{problem}

\section{Questions d'entretien}

\begin{interviewq}[$\star$ \iqroles{trader}]\label{iq:fm:treasury-and-funding:1}
Qu'est-ce que la \omterm{def:fm:treasury-and-funding:house}{marge maison} (house margin), et pourquoi est-elle généralement supérieure au minimum réglementaire~?
\end{interviewq}

\begin{interviewq}[$\star$ \iqroles{risque}]\label{iq:fm:treasury-and-funding:2}
Qu'est-ce qu'un \omterm{def:fm:treasury-and-funding:horizon}{horizon de survie}, et comment le calculeriez-vous~?
\end{interviewq}

\begin{interviewq}[$\star\star$ \iqroles{chercheur, développeur}]\label{iq:fm:treasury-and-funding:3}
Formulez l'allocation des positions entre courtiers, avec les termes de marge brute, nette et de concentration, comme un problème d'optimisation.
\end{interviewq}

\begin{interviewq}[$\star\star$ \iqroles{risque}]\label{iq:fm:treasury-and-funding:4}
Votre prime broker double sa marge pendant la nuit. Que vérifiez-vous, et que faites-vous~?
\end{interviewq}

\begin{interviewq}[$\star\star$ \iqroles{trader, risque}]\label{iq:fm:treasury-and-funding:5}
Pourquoi une société pourrait-elle payer plus de marge qu'il n'en faut~?
\end{interviewq}

\begin{interviewq}[$\star\star\star$ \iqroles{risque, chercheur}]\label{iq:fm:treasury-and-funding:6}
Concevez un test de résistance de liquidité pour un hedge fund disposant de trois prime brokers et d'un compte de futures.
\end{interviewq}
